\documentclass[10pt,a4paper]{amsart}
\usepackage[margin=19mm]{geometry}
\usepackage{amsmath,amssymb,mathtools}
\usepackage[scr=rsfs,cal=euler]{mathalfa}
\usepackage{xfrac}
\usepackage[unicode,hidelinks,bookmarks=false]{hyperref}
\newcommand{\ie}{i.e.\ }
\newcommand{\cf}{cf.\ }
\newcommand{\resp}{respectively\ }
\newcommand{\Subsection}{\subsection}
\providecommand{\nobreakdash}{\nobreak\hbox{-}}
\setlength{\emergencystretch}{2em}
\multlinegap0pt
\usepackage{amssymb}
\newtheorem{lemma}{Lemma}
\title{Extremal problem of area approximation}
\author{O.O. Pokutnyi, R.R. Salimov, M.V. Stefanchuk}
\date{Source: arXiv:2607.24828v2 (CC BY 4.0)}
\begin{document}
\maketitle

\noindent\emph{Source and licence.} Extremal problem of area approximation. Version arXiv:2607.24828v2 (CC BY 4.0), distributed by arXiv under the Creative Commons Attribution 4.0 International licence, http://creativecommons.org/licenses/by/4.0/, as recorded in the arXiv licence metadata for this exact version and shown on the abstract page https://arxiv.org/abs/2607.24828v2. Full licence notice: \texttt{licenses/arxiv-2607.24828-CC-BY-4.0.txt}.\par\smallskip

\noindent\textit{Editorial scope.} Counted unit: the article's lemma lemma4 (source0.tex lines 203--224). The article's complete original proof of it is delivered with it (source0.tex lines 227--278, one \texttt{proof} environment of 52 lines). No mathematics is written, repaired or re-derived: the statement and the proof are the article's own lines, in the article's own notation, and no retained line is substituted or re-flowed.  Retained uncounted as the source-local context the proof needs: lines 179--184.  Nothing was omitted from any retained span, so the package carries no omission record.  No retained line contains a citation, so the wrapper carries no bibliography and no bibliography entry.  No retained line contains a figure environment, an image-inclusion command or drawing code, so no image is delivered and the package declares no asset dependency.  Editorial formatting only, in addition: an AMS \texttt{amsart} preamble that restates the article's own declarations and macros used by the retained lines, namely the environments \texttt{lemma} on separate counters, together with the standard packages those lines need.  The retained environments print in this document as Lemma 1 in order of appearance, whereas the article prints the same items with the numbering of its own full document.  The complete native source of the article as distributed in the arXiv e-print for this exact version is bundled unchanged as \texttt{sources/206144/source0.tex}.
\begin{lemma}\label{lemma3}
Let $f(x) = Ax^2 + Bx + C$, $A \neq 0$, and let $a =x_0 < x_1 < \cdots < x_n < x_{n+1} = b$ be a partition. Let $\ell_i(x)$ be the chord connecting the points $(x_{i-1}, f(x_{i-1}))$ and $(x_i, f(x_i))$. Then the total area $S(x_1, \ldots, x_n)$ between the graph of $f(x)$ and the polygonal line formed by these chords is given by
\[
S(x_1, \ldots, x_n) = \frac{|A|}{6} \sum_{i=1}^{n+1} (x_i - x_{i-1})^3.
\]
\end{lemma}

\begin{lemma}\label{lemma4}
Let
\[
S(x_1, \dots, x_n) = \frac{|A|}{6} \sum_{i=1}^{n+1} (x_i - x_{i-1})^3, \qquad x_0 = a, \quad x_{n+1} = b, \quad A\neq 0.
\]
Then the Hessian matrix
$
H = \left( \frac{\partial^2 S}{\partial x_i \partial x_j} \right)_{i,j=1}^{n}
$
of size \( n \times n \) has the following form:
\[
H = |A| \cdot
\begin{pmatrix}
x_2 - x_0 & -(x_2 - x_1) & 0 & \cdots & 0 \\
-(x_2 - x_1) & x_3 - x_1 & -(x_3 - x_2) & \cdots & 0 \\
0 & -(x_3 - x_2) & x_4 - x_2 & \cdots & 0 \\
\vdots & \vdots & \vdots & \ddots & \vdots \\
0 & 0 & 0 & \cdots & x_{n+1} - x_{n-1}
\end{pmatrix}.
\]

\end{lemma}

\begin{proof}
By Lemma \ref{lemma3}, we have
\[
S(x_1, \dots, x_n) = \frac{|A|}{6} \sum_{k=1}^{n+1} (x_k - x_{k-1})^3.
\]
Only two terms in this sum depend on \( x_i \) (those with indices \( k = i \) and \( k = i+1 \)):
\[
S = \frac{|A|}{6} \Bigl( (x_i - x_{i-1})^3 + (x_{i+1} - x_i)^3 \Bigr) + \left( S-\frac{|A|}{6} \Bigl( (x_i - x_{i-1})^3 + (x_{i+1} - x_i)^3 \Bigr)\right).
\]
Compute the first derivative with respect to $x_i$ of the function $S$:
\[
\frac{\partial S}{\partial x_i}
= \frac{|A|}{2} \Bigl( (x_i - x_{i-1})^2 - (x_{i+1} - x_i)^2 \Bigr).
\]
Next, we find the second derivatives of $S$.

Diagonal entry:
\[
\frac{\partial^2 S}{\partial x_i^2}
= |A| \bigl(  x_{i+1} - x_{i-1}  \bigr).
\]

Off-diagonal entry \( (i, i+1) \):
\[
\frac{\partial^2 S}{\partial x_i \partial x_{i+1}} = \frac{\partial}{\partial x_{i+1}} \left( \frac{|A|}{2} \bigl( (x_i - x_{i-1})^2 - (x_{i+1} - x_i)^2 \bigr) \right)
= -|A| (x_{i+1} - x_i).
\]

Entries \( (i, i-1) \): by symmetry of the Hessian matrix,
\[
\frac{\partial^2 S}{\partial x_i \partial x_{i-1}} = \frac{\partial^2 S}{\partial x_{i-1} \partial x_i} = -|A| (x_i-x_{i-1}).
\]

All other entries: if \( |i - j| > 1 \), then \( x_i \) and \( x_j \) appear in different terms of the sum, hence,
\[
\frac{\partial^2 S}{\partial x_i \partial x_j} = 0.
\]

Thus, the Hessian matrix has the form
\[
H = |A| \cdot
\begin{pmatrix}
x_2 - x_0 & -(x_2 - x_1) & 0 & \cdots & 0 \\
-(x_2 - x_1) & x_3 - x_1 & -(x_3 - x_2) & \cdots & 0 \\
0 & -(x_3 - x_2) & x_4 - x_2 & \cdots & 0 \\
\vdots & \vdots & \vdots & \ddots & \vdots \\
0 & 0 & 0 & \cdots & x_{n+1} - x_{n-1}
\end{pmatrix}.
\]

Lemma \ref{lemma4} is proved. 
\end{proof}

\end{document}
